\begin{figure}[!t]
\centering
\colorlet{cadv}{red!80!black}
% Blocks sit 0.52 apart in height, k = 4. The trunk runs from the genesis
% block (height 0) to height 9. Two honest chains leave it: one ends at height
% 13, the other at height 12. Two adversarial chains: one leaves after the
% second block (height 1) and has 7 blocks (heights 2..8), the other leaves
% after height 4 and has 5 blocks (heights 5..9). Parties stand right of tips.
\begin{tikzpicture}[font=\footnotesize]
  \tikzset{adv/.style={blk, draw=cadv}, advcut/.style={cut, draw=cadv}}
  \def\yU{0.8}    % upper honest chain
  \def\yA{-1.3}   % lower adversarial chain
  \def\yB{1.75}   % upper adversarial chain

  % height axis and the estimate of the client
  \draw[->] (-0.3,-2.35) -- (8.2,-2.35) node[right] {height};
  \draw[lim] (7.6,-2.35) -- (7.6,2.75)
    node[above,align=center] {honest length\\at round $r$};

  % trunk, the common prefix: heights 0..9. Height 9 is among the last k
  % blocks of the shorter honest chain, so it is shaded.
  \foreach \i in {0,...,8} { \blk{0.52*\i}{0} }
  \blk[cut]{4.68}{0}
  \foreach \i in {0,...,8} { \lnk{0.52*\i}{0}{0.52*\i+0.52}{0} }
  \node[below=3pt] at (0,-0.21) {genesis};

  % longer honest chain: heights 10..13 on the trunk line
  \foreach \i in {10,...,13} { \blk[cut]{0.52*\i}{0} }
  \lnk{4.68}{0}{5.2}{0}
  \foreach \i in {10,...,12} { \lnk{0.52*\i}{0}{0.52*\i+0.52}{0} }
  \pic at (7.18,-0.3) {person};
  \draw[brc] (5.03,-0.26) -- (6.93,-0.26) node[midway,below=4pt] {$k$};

  % shorter honest chain: heights 10..12
  \foreach \i in {10,...,12} { \blk[cut]{0.52*\i}{\yU} }
  \lnk{4.68}{0}{5.2}{\yU}
  \foreach \i in {10,...,11} { \lnk{0.52*\i}{\yU}{0.52*\i+0.52}{\yU} }
  \pic at (6.66,\yU-0.3) {person};

  % first adversarial chain: leaves after height 1, heights 2..8
  \draw[ink,cadv] (0.69,0) -- (0.87,\yA);
  \foreach \i in {2,...,4} { \blk[adv]{0.52*\i}{\yA} }
  \foreach \i in {5,...,8} { \blk[advcut]{0.52*\i}{\yA} }
  \foreach \i in {2,...,7} { \draw[ink,cadv] (0.52*\i+0.17,\yA) -- (0.52*\i+0.35,\yA); }
  \pic[cadv] at (4.58,\yA-0.3) {devil};
  \draw[brc] (0.87,\yA-0.26) -- (4.33,\yA-0.26)
    node[midway,below=4pt] {more than $k$};

  % second adversarial chain: leaves after height 4, heights 5..9
  \draw[ink,cadv] (2.25,0) -- (2.43,\yB);
  \blk[adv]{2.6}{\yB}
  \foreach \i in {6,...,9} { \blk[advcut]{0.52*\i}{\yB} }
  \foreach \i in {5,...,8} { \draw[ink,cadv] (0.52*\i+0.17,\yB) -- (0.52*\i+0.35,\yB); }
  \pic[cadv] at (5.1,\yB-0.3) {devil};
  \draw[brc] (4.85,\yB+0.26) -- (2.43,\yB+0.26)
    node[midway,above=4pt] {more than $k$};
\end{tikzpicture}
\caption{The block tree at round $r$, drawn by height. The dashed line is the
height that the client estimates, from its clock, for the chains of the honest
validators. Two honest validators hold chains of different lengths that end
next to it and share every block up to a short fork. The last $k$ blocks of
each are shaded. The adversary, in red, holds two much shorter chains. Each
leaves the chain of the honest validators more than $k$ blocks below its tip.
Their last $k$ blocks are shaded too, and the rest of each lies off the chain
of the honest validators.}\label{fig.tree}
\end{figure}
